% !TeX root = ../Book1.tex
\begin{figure}[htbp]
	\centering
	\begin{tikzpicture}[font=\small]
		\fill[GeometryBlueLight] (0,0) rectangle (6,2.4);
		\draw[GeometryGray,dashed] (0,0)--(0,2.4);
		\draw[GeometryGray,dashed] (6,0)--(6,2.4);
		\draw[GeometryBlue,thick,->] (0,0)--(3.6,0);
		\draw[GeometryBlue,thick] (3.6,0)--(6,0);
		\draw[GeometryBlue,thick,->] (0,2.4)--(3.6,2.4);
		\draw[GeometryBlue,thick] (3.6,2.4)--(6,2.4);
		\draw[GeometryGray,->] (0,1.2)--(6,1.2);
		\draw[orange!80!black,thick,->] (2.6,0.4)--(2.6,1.95);
		\node[right,orange!80!black] at (2.6,1.75) {改变路径：\(\partial_s\)};
		\node[below] at (3,0) {沿路径前进：\(\partial_t\)};
		\node[above] at (3,2.4) {\(\Gamma(s_1,t)\)};
		\node[left,align=right] at (-0.1,1.2) {起点曲线\\\(\Gamma(s,a)\)};
		\node[right,align=left] at (6.1,1.2) {终点曲线\\\(\Gamma(s,b)\)};
		\node[below left] at (0,0) {\((s_0,a)\)};
		\node[below right] at (6,0) {\((s_0,b)\)};
	\end{tikzpicture}
	\caption{图画在路径族的参数域中，\(\Gamma\) 本身可以自交或降秩。比较相邻路径的运输时，中间窄带给出曲率项；两端的移动给出端点项。只有定端点变化才能直接省去这两项。}
	\label{g1:fig:principal-transport-variation}
\end{figure}
